\subsubsection{Kiểm tra giả thiết Gauss--Markov}
\paragraph{Biểu đồ phần dư (residual plot)}

\figurename~\ref{fig:ols-residual-plot} thể hiện phần dư trên tập huấn luyện.
Normal Equation được dùng để tìm nghiệm bình phương tối thiểu. Nhìn vào hình ta
thấy có xuất hiện vệt của đường thẳng. Đây là do việc ngắt ngưỡng của biến mục
tiểu. Độ phân tán không đồng đều (trong đoạn $[0,1]$ và $[4, 6]$) thấp và không
ổn định như trong đoạn $[2, 4]$. Điều này cho thấy giả thiết về phương sai bằng
nhau giữa các mẫu là không đúng. Phân tán phần dư trong đoạn $[0, 1]$ và $[4,
6]$ nằm trên đường thẳng $e = 0$ cho thấy trung bình phần dư trên đoạn này cũng
không bằng 0. Trên đoạn $[6, 8]$ cũng thấy trung bình phần dư nhỏ hơn 0 nhưng
nguyên nhân chính là do việc ngắt ngưỡng. Trên đoạn $[2, 4]$, giả thiết
Gauss--Markov vẫn giữ tính đúng.

\begin{figure}[htbp]
  \centering
  \includegraphics[width=0.9\textwidth]{ols-residual-plot}
  \caption{Biểu đồ phần dư của mô hình hồi quy tuyến tính qua phương pháp bình
  phương tối thiểu}
  \label{fig:ols-residual-plot}
\end{figure}

\paragraph{Biểu đồ phân vị-phân vị (Quantile-Quantile plot)}

\figurename~\ref{fig:ols-qq-plot} cho thấy sự tuân theo phân phối chuẩn của
phần dư trên tập huấn luyện. Ở vùng trung tâm ($[-1, 2]$) các điểm bám sát với
đường tham chiếu. Điều này thể hiện mô hình ổn định và phần dư bám sát phân
phối chuẩn. Vùng đuôi bên trái ($[-3, -11]$), đường biểu diễn cong nhẹ lên đường
tham chiếu. Điều này cho thấy hiện tượng đuôi mỏng, nghĩa là mô hình hiếm khi
xảy ra sai số âm cực lớn. Vùng đuôi bên phải ($[2, 4]$), xuất hiện hiện tượng
vểnh cao rõ rệt. Đây là dấu hiệu hiện tượng đuôi dày, phản ánh sự tồn tại của
các sai số dương lớn. Hiện tượng này chủ yếu do dữ liệu bị chặn ngưỡng tại mức.
Khi giá trị thực tế cao vượt ngưỡng nhưng dữ liệu ghi nhận thấp hơn, mô hình
tuyến tính OLS khó có thể khớp chính xác, dẫn đến bùng nổ sai số dương ở phân
khúc này.

\begin{figure}[htbp]
  \centering
  \includegraphics[width=0.9\textwidth]{ols-qq-plot}
  \caption{Biểu đồ phân vị-phân vị phần dư của mô hình hồi quy tuyến tính
  qua phương pháp bình phương tối thiểu khi so với phân
  phối chuẩn}
  \label{fig:ols-qq-plot}
\end{figure}

\paragraph{Kiểm định Breusch--Pagan}

\tablename~\ref{tab:breusch-pagan-test-result} cho kết quả đánh giá giả thuyết
về phương sai sai số không đổi (Homoscedasticity)---một trong điều kiện tiên
quyết của định lý Gauss--Markov để ước lượng bình phương tối thiểu là không
chệch tuyến tính tốt nhất (BLUE). Cả hai giá trị thống kê (LM Stat khoảng
$1248.62$ và F Stat khoảng $173.45$) đều rất lớn. Điều này cho thấy sự sai lệch
đáng kê giữa phân phối phần dư thực tế và giả thiết phương sai sai số đồng
nhất. Và với mức ý nghĩa $\alpha = 0.05$, ta có $p$-value nhỏ hơn $\alpha$. Do
đó ta bác bỏ giả thuyết $H_0$ (phương sai sai số không đổi) và chấp nhận giả
thuyết $H_1$ (phương sai sai số thay đổi). Kết quả này nhất quán với nội dung
biểu đồ phần dư khi độ rộng phân tán không đều.

\begin{table}[htbp]
  \centering
  \caption{Kết quả kiểm định Breusch--Pagan trên tập huấn luyện}
  \label{tab:breusch-pagan-test-result}
  \import{\tablespath}{breusch-pagan-test-result}
\end{table}
